% ============================================================================
% EMCH 792: Optimal State Estimation - Homework 1 (Written Portion)
% ============================================================================

\documentclass[12pt]{article}

% ============================================================================
% PACKAGE IMPORTS
% ============================================================================
\usepackage{amsmath}
\usepackage{graphicx}
\usepackage{amssymb}
\usepackage{tikz}
\usepackage[margin=1in]{geometry}
\usepackage{setspace}
\usepackage{xcolor}
\usepackage{enumitem}
\usepackage{tcolorbox}
\usepackage{fancyhdr}
\usepackage{booktabs}
\usepackage{array}
\usepackage{colortbl}
%\usepackage{pgfplots}
\usepackage{listings}
\usepackage{float}
%\pgfplotsset{compat=1.18}

% ============================================================================
% HEADER AND FOOTER SETUP
% ============================================================================
\pagestyle{fancy}
\setlength{\headheight}{14.5pt}
\fancyhf{}

\fancyhead[L]{Page \thepage}
\fancyhead[C]{EMCH 792}
\fancyhead[R]{JC Vaught}

\renewcommand{\headrulewidth}{0pt}

% ============================================================================
% TIKZ LIBRARIES
% ============================================================================
\usetikzlibrary{shapes.geometric, arrows.meta, positioning, decorations.pathmorphing, patterns}

% --- USC Brand Colors ---
\definecolor{USC_Garnet}{HTML}{73000A}
\definecolor{USC_Sandstorm}{HTML}{FFF2E3}
\definecolor{USC_Black90}{HTML}{363636}
\definecolor{USC_Black10}{HTML}{ECECEC}
\definecolor{USC_Honeycomb}{HTML}{A49137}
\definecolor{tablegreen}{RGB}{230,242,230}
\definecolor{outputbg}{RGB}{245,245,245}

% ============================================================================
% CUSTOM COLORED BOXES
% ============================================================================
\newtcolorbox{stepbox}{
  colback=white,
  colframe=USC_Garnet,
  fonttitle=\bfseries,
  title=Step,
  sharp corners,
  colbacktitle=USC_Garnet,
  coltitle=white,
}

\newtcolorbox{codebox}{
  colback=USC_Black10,
  colframe=USC_Black90,
  fonttitle=\bfseries,
  title=Python Implementation,
  sharp corners,
  colbacktitle=USC_Black90,
  coltitle=white,
}

\newtcolorbox{outputbox}{
  colback=outputbg,
  colframe=USC_Honeycomb,
  fonttitle=\bfseries,
  title=Python Output,
  sharp corners,
  colbacktitle=USC_Honeycomb,
  coltitle=white,
  breakable,
}

\newtcolorbox{resultsbox}{
  colback=white,
  colframe=USC_Honeycomb,
  fonttitle=\bfseries,
  title=Results,
  sharp corners,
  colbacktitle=USC_Honeycomb,
  coltitle=white,
}

% ============================================================================
% CUSTOM QUESTION AND PART COMMANDS
% ============================================================================
\newcommand{\question}[1]{%
  \clearpage
  \vspace{0.5cm}
  {\noindent\normalsize \textbf{#1}}
  \vspace{0.2cm}
  \hrule
  \vspace{0.1cm}
  \hrule
  \vspace{0.3cm}
}

\newcommand{\questionpart}[1]{%
  \clearpage
  \vspace{0.5cm}
  {\noindent\normalsize \textbf{#1}}
  \vspace{0.2cm}
  \hrule
  \vspace{0.1cm}
  \hrule
  \vspace{0.3cm}
}

% Code listing style
\lstset{
    language=Python,
    basicstyle=\ttfamily\small,
    keywordstyle=\color{blue},
    commentstyle=\color{USC_Honeycomb},
    stringstyle=\color{USC_Garnet},
    breaklines=true,
    showstringspaces=false,
    backgroundcolor=\color{USC_Black10},
    numbers=left,
    numberstyle=\tiny\color{gray},
    numbersep=5pt,
}

% Output listing style
\lstdefinestyle{output}{
    basicstyle=\ttfamily\footnotesize,
    breaklines=true,
    showstringspaces=false,
    backgroundcolor=\color{outputbg},
    frame=none,
    numbers=none,
}

% Graphics path
\graphicspath{{figures/}}

% ============================================================================
% DOCUMENT INFORMATION
% ============================================================================

\title{EMCH 792: Optimal State Estimation \\ Homework 1 \\ \large{Written Portion Solutions}}
\author{Instructor: N. Vitzilaios, Department of Mechanical Engineering \\ University of South Carolina \\ \textit{Solutions by: JC Vaught}}
\date{Due: January 28, 2026}

% ============================================================================
% BEGIN DOCUMENT
% ============================================================================

\begin{document}

\maketitle
\setlength{\parindent}{0pt}

\setlist[enumerate,1]{label=\arabic*.}
\setlist[enumerate,2]{label=(\alph*)}

% ============================================================================
% TABLE OF CONTENTS
% ============================================================================

\begin{center}
\begin{tcolorbox}[colback=white, colframe=gray!50!black, colbacktitle=gray!20!white,
                  coltitle=black, sharp corners, boxrule=1pt,
                  title=\Large\bfseries Table of Contents]
\begin{tabular}{p{0.82\textwidth}r}
\textbf{Exercise 1.2:} Commuting Non-Diagonal Matrices (10 pts) \dotfill & \textbf{\pageref{quest:1.2}} \\[0.3em]
\textbf{Exercise 1.3:} Prove Identities of Equation (1.26) (10 pts) \dotfill & \textbf{\pageref{quest:1.3}} \\[0.3em]
\textbf{Exercise 1.4:} Partial Derivative of $\mathrm{Tr}(AB)$ w.r.t.\ $A$ (10 pts) \dotfill & \textbf{\pageref{quest:1.4}} \\[0.3em]
\textbf{Exercise 1.6:} Block Matrix Inverse (10 pts) \dotfill & \textbf{\pageref{quest:1.6}} \\[0.3em]
\textbf{Exercise 1.7:} Symmetry of $AB$ (10 pts) \dotfill & \textbf{\pageref{quest:1.7}} \\
\end{tabular}
\vspace{0.2cm}
\end{tcolorbox}
\end{center}

\vspace{0.5cm}

% ============================================================================
% EXERCISE 1.2
% ============================================================================

\question{Exercise 1.2: Commuting Non-Diagonal Matrices (10 pts)}\label{quest:1.2}

Find two $2 \times 2$ matrices $A$ and $B$ such that $A \neq B$, neither $A$ nor $B$ are diagonal, $A \neq cB$ for any scalar $c$, and $AB = BA$. Find the eigenvectors of $A$ and $B$.

\begin{stepbox}
\textbf{Construction}

We need two non-diagonal, non-scalar-multiple matrices that commute. A standard approach: choose matrices that share the same eigenvectors but have different eigenvalues.

Let
\[
A = \begin{pmatrix} 1 & 1 \\ 0 & 1 \end{pmatrix}, \qquad
B = \begin{pmatrix} 1 & 2 \\ 0 & 1 \end{pmatrix}
\]

\textbf{Verification of conditions:}
\begin{itemize}
    \item $A \neq B$ since $a_{12} = 1 \neq 2 = b_{12}$. \checkmark
    \item Neither is diagonal. \checkmark
    \item $A \neq cB$: if $A = cB$ then $c \cdot 1 = 1$ gives $c=1$, but $c \cdot 2 = 1$ gives $c = 1/2$. Contradiction. \checkmark
\end{itemize}
\end{stepbox}

\begin{stepbox}
\textbf{Checking $AB = BA$}

\[
AB = \begin{pmatrix} 1 & 1 \\ 0 & 1 \end{pmatrix}\begin{pmatrix} 1 & 2 \\ 0 & 1 \end{pmatrix} = \begin{pmatrix} 1 & 3 \\ 0 & 1 \end{pmatrix}
\]

\[
BA = \begin{pmatrix} 1 & 2 \\ 0 & 1 \end{pmatrix}\begin{pmatrix} 1 & 1 \\ 0 & 1 \end{pmatrix} = \begin{pmatrix} 1 & 3 \\ 0 & 1 \end{pmatrix}
\]

$AB = BA$. \checkmark
\end{stepbox}

\begin{stepbox}
\textbf{Eigenvectors}

Both $A$ and $B$ are upper triangular with repeated eigenvalue $\lambda = 1$.

For $A$: $(A - I)v = 0$ gives $\begin{pmatrix} 0 & 1 \\ 0 & 0 \end{pmatrix}v = 0$, so $v_2 = 0$. Eigenvector: $\mathbf{v} = \begin{pmatrix} 1 \\ 0 \end{pmatrix}$.

For $B$: $(B - I)v = 0$ gives $\begin{pmatrix} 0 & 2 \\ 0 & 0 \end{pmatrix}v = 0$, so $v_2 = 0$. Eigenvector: $\mathbf{v} = \begin{pmatrix} 1 \\ 0 \end{pmatrix}$.
\end{stepbox}

\begin{resultsbox}
\[
A = \begin{pmatrix} 1 & 1 \\ 0 & 1 \end{pmatrix}, \quad
B = \begin{pmatrix} 1 & 2 \\ 0 & 1 \end{pmatrix}, \quad
AB = BA = \begin{pmatrix} 1 & 3 \\ 0 & 1 \end{pmatrix}
\]
Both $A$ and $B$ share the eigenvector $\boxed{\mathbf{v} = \begin{pmatrix} 1 \\ 0 \end{pmatrix}}$ (with eigenvalue $\lambda = 1$).
\end{resultsbox}

% ============================================================================
% EXERCISE 1.3
% ============================================================================

\question{Exercise 1.3: Prove the Three Identities of Equation (1.26) (10 pts)}\label{quest:1.3}

Prove:
\begin{enumerate}
    \item $(AB)^T = B^T A^T$
    \item $(AB)^{-1} = B^{-1} A^{-1}$
    \item $\mathrm{Tr}(AB) = \mathrm{Tr}(BA)$
\end{enumerate}

\begin{stepbox}
\textbf{Proof of $(AB)^T = B^T A^T$}

Let $A$ be $m \times n$ and $B$ be $n \times p$. Then $AB$ is $m \times p$.

The $(i,j)$ entry of $(AB)^T$ is the $(j,i)$ entry of $AB$:
\[
[(AB)^T]_{ij} = (AB)_{ji} = \sum_{k=1}^{n} a_{jk}\, b_{ki}
\]

The $(i,j)$ entry of $B^T A^T$:
\[
[B^T A^T]_{ij} = \sum_{k=1}^{n} (B^T)_{ik}\, (A^T)_{kj} = \sum_{k=1}^{n} b_{ki}\, a_{jk} = \sum_{k=1}^{n} a_{jk}\, b_{ki}
\]

These are identical, so $(AB)^T = B^T A^T$. \hfill $\blacksquare$
\end{stepbox}

\begin{stepbox}
\textbf{Proof of $(AB)^{-1} = B^{-1} A^{-1}$}

Assume $A$ and $B$ are square and invertible. We verify that $B^{-1}A^{-1}$ is the inverse of $AB$ by showing it satisfies both conditions:

\textbf{Right inverse:}
\[
(AB)(B^{-1}A^{-1}) = A(BB^{-1})A^{-1} = AIA^{-1} = AA^{-1} = I
\]

\textbf{Left inverse:}
\[
(B^{-1}A^{-1})(AB) = B^{-1}(A^{-1}A)B = B^{-1}IB = B^{-1}B = I
\]

Since the inverse is unique, $(AB)^{-1} = B^{-1}A^{-1}$. \hfill $\blacksquare$
\end{stepbox}

\begin{stepbox}
\textbf{Proof of $\mathrm{Tr}(AB) = \mathrm{Tr}(BA)$}

Let $A$ be $m \times n$ and $B$ be $n \times m$ (so both $AB$ and $BA$ are square).

\[
\mathrm{Tr}(AB) = \sum_{i=1}^{m} (AB)_{ii} = \sum_{i=1}^{m} \sum_{k=1}^{n} a_{ik}\, b_{ki}
\]

\[
\mathrm{Tr}(BA) = \sum_{k=1}^{n} (BA)_{kk} = \sum_{k=1}^{n} \sum_{i=1}^{m} b_{ki}\, a_{ik}
\]

Both double sums range over the same indices and contain the same terms $a_{ik}\, b_{ki}$, just with the order of summation swapped. Therefore $\mathrm{Tr}(AB) = \mathrm{Tr}(BA)$. \hfill $\blacksquare$
\end{stepbox}

% ============================================================================
% EXERCISE 1.4
% ============================================================================

\question{Exercise 1.4: Partial Derivative of $\mathrm{Tr}(AB)$ with Respect to $A$ (10 pts)}\label{quest:1.4}

Find $\dfrac{\partial}{\partial A}\,\mathrm{Tr}(AB)$.

\begin{stepbox}
\textbf{Derivation}

Let $A$ be $m \times n$ and $B$ be $n \times m$. Then:
\[
\mathrm{Tr}(AB) = \sum_{i=1}^{m} (AB)_{ii} = \sum_{i=1}^{m} \sum_{k=1}^{n} a_{ik}\, b_{ki}
\]

The matrix derivative $\dfrac{\partial}{\partial A}\,\mathrm{Tr}(AB)$ is defined as the $m \times n$ matrix whose $(i,j)$ entry is:
\[
\left[\frac{\partial}{\partial A}\,\mathrm{Tr}(AB)\right]_{ij} = \frac{\partial}{\partial a_{ij}}\,\mathrm{Tr}(AB)
\]

Since $\mathrm{Tr}(AB) = \sum_{p=1}^{m}\sum_{k=1}^{n} a_{pk}\,b_{kp}$, the only term involving $a_{ij}$ is $a_{ij}\,b_{ji}$. Therefore:
\[
\frac{\partial}{\partial a_{ij}}\,\mathrm{Tr}(AB) = b_{ji}
\]

This means $\left[\dfrac{\partial}{\partial A}\,\mathrm{Tr}(AB)\right]_{ij} = b_{ji} = (B^T)_{ij}$.
\end{stepbox}

\begin{resultsbox}
\[
\boxed{\frac{\partial}{\partial A}\,\mathrm{Tr}(AB) = B^T}
\]
\end{resultsbox}

% ============================================================================
% EXERCISE 1.6
% ============================================================================

\question{Exercise 1.6: Block Matrix Inverse (10 pts)}\label{quest:1.6}

Suppose that $A$ is invertible and
\[
\begin{bmatrix} A & A \\ B & A \end{bmatrix}
\begin{bmatrix} A \\ C \end{bmatrix}
=
\begin{bmatrix} 0 \\ I \end{bmatrix}
\]
Find $B$ and $C$ in terms of $A$.

\begin{stepbox}
\textbf{Expanding the Block Multiplication}

Performing the block matrix--vector multiplication:
\[
\begin{bmatrix} A \cdot A + A \cdot C \\ B \cdot A + A \cdot C \end{bmatrix}
=
\begin{bmatrix} 0 \\ I \end{bmatrix}
\]

This gives two equations:
\begin{align}
A^2 + AC &= 0 \tag{i}\\
BA + AC &= I \tag{ii}
\end{align}
\end{stepbox}

\begin{stepbox}
\textbf{Solving for $C$}

From equation (i):
\[
A^2 + AC = 0 \implies A(A + C) = 0
\]

Since $A$ is invertible, we can left-multiply by $A^{-1}$:
\[
A + C = 0 \implies \boxed{C = -A}
\]
\end{stepbox}

\begin{stepbox}
\textbf{Solving for $B$}

Substitute $C = -A$ into equation (ii):
\[
BA + A(-A) = I \implies BA - A^2 = I
\]
\[
BA = I + A^2
\]

Right-multiply by $A^{-1}$:
\[
B = (I + A^2)A^{-1} = A^{-1} + A
\]
\end{stepbox}

\begin{resultsbox}
\[
\boxed{C = -A}, \qquad \boxed{B = A^{-1} + A}
\]

\textbf{Verification:} Substituting back into (i): $A^2 + A(-A) = A^2 - A^2 = 0$ \checkmark

Substituting back into (ii): $(A^{-1} + A)A + A(-A) = I + A^2 - A^2 = I$ \checkmark
\end{resultsbox}

% ============================================================================
% EXERCISE 1.7
% ============================================================================

\question{Exercise 1.7: $AB$ May Not Be Symmetric (10 pts)}\label{quest:1.7}

Show that $AB$ may not be symmetric even though both $A$ and $B$ are symmetric.

\begin{stepbox}
\textbf{Counterexample}

Let
\[
A = \begin{pmatrix} 1 & 2 \\ 2 & 3 \end{pmatrix}, \qquad
B = \begin{pmatrix} 1 & 0 \\ 0 & 2 \end{pmatrix}
\]

Both are symmetric: $A^T = A$ and $B^T = B$.

\[
AB = \begin{pmatrix} 1 & 2 \\ 2 & 3 \end{pmatrix}\begin{pmatrix} 1 & 0 \\ 0 & 2 \end{pmatrix} = \begin{pmatrix} 1 & 4 \\ 2 & 6 \end{pmatrix}
\]

Check symmetry:
\[
(AB)^T = \begin{pmatrix} 1 & 2 \\ 4 & 6 \end{pmatrix} \neq \begin{pmatrix} 1 & 4 \\ 2 & 6 \end{pmatrix} = AB
\]

Since $(AB)^T \neq AB$, the product $AB$ is not symmetric. \hfill $\blacksquare$
\end{stepbox}

\begin{stepbox}
\textbf{General Explanation}

For $AB$ to be symmetric, we would need $(AB)^T = AB$, i.e., $B^T A^T = AB$. When $A$ and $B$ are symmetric this becomes $BA = AB$. So $AB$ is symmetric if and only if $A$ and $B$ commute. In general, symmetric matrices do not commute, so their product need not be symmetric.
\end{stepbox}


% ============================================================================
% END OF DOCUMENT
% ============================================================================

\end{document}
